% 文件开头添加条件判断
\newif\ifstandalone
\standalonetrue  % 默认独立编译模式

% 检测是否在主文档中
\ifdefined\maindoc
  \standalonefalse
\fi

\ifstandalone
  % 第二章：控制理论导论
  % Chapter 2: Introduction to Control Theory
  \documentclass[UTF8,12pt,a4paper]{ctexart}

  % 页面设置
  \usepackage{geometry}
  \geometry{left=2.5cm,right=2.5cm,top=2.5cm,bottom=2.5cm}

  % 数学包
  \usepackage{amsmath}
  \usepackage{amssymb}
  \usepackage{amsthm}

  % 图形和表格
  \usepackage{graphicx}
  \usepackage{array}
  \usepackage{booktabs}

  % 其他包
  \usepackage{hyperref}
  \hypersetup{colorlinks=true,linkcolor=blue,citecolor=blue}

  % 定理环境定义
  \theoremstyle{definition}
  \newtheorem{definition}{定义}[section]
  \newtheorem{theorem}{定理}[section]
  \newtheorem{lemma}{引理}[section]
  \newtheorem{example}{示例}[section]
  \newtheorem{remark}{注记}[section]

  % 章节标题设置
  \ctexset{
    section = {
      format = \Large\bfseries\raggedright,
      beforeskip = 1.5ex plus 0.5ex minus 0.5ex,
      afterskip = 1.0ex plus 0.2ex minus 0.2ex
    },
    subsection = {
      format = \large\bfseries\raggedright,
      beforeskip = 1.0ex plus 0.3ex minus 0.3ex,
      afterskip = 0.5ex plus 0.1ex minus 0.1ex
    }
  }
\fi

\ifstandalone
  \begin{document}
  % 章节标题
  \section*{第二章：控制理论导论}
  \addcontentsline{toc}{section}{第二章：控制理论导论}
\fi

% 当在主文档中使用时，只保留章节内容
\ifstandalone
\else
  \section*{第二章：控制理论导论}
  \addcontentsline{toc}{section}{第二章：控制理论导论}
\fi

在本章中，我们回顾了用于机器人控制的各种控制理论概念。我们首先回顾了线性和非线性系统的状态空间描述，并介绍了后续所需的稳定性概念。本章旨在介绍现代控制概念，但即使是具有控制理论背景的读者也可能希望查阅它以了解符号和方便性。

\section{引言}

机器人机械臂的控制是一个成熟但富有成效的研究、开发和制造领域。工业机器人本质上是定位和搬运设备。因此，一个有用的机器人是能够控制其运动以及机器人与环境之间的相互作用力和力矩的机器人。本书关注机器人机械臂的控制方面。控制通常需要数学模型的可用性以及某种作用于模型的智能。机器人的数学模型是从控制其运动的基本物理定律中获得的。另一方面，智能需要感知能力以及作用于和响应感知变量的手段。这些机器人的作用和反应是控制器设计的结果。

在本章中，我们回顾了本书所需的控制理论概念。所有证明都被省略，但参考了提供更多证明的更专业书籍和论文。一旦在第2章中描述了机器人动力学的满意模型，本章介绍的自动控制概念可用于修改机器人对不同刺激的作用和反应。因此，后续章节将讨论将控制原理应用于机器人方程。所使用的特定控制器将取决于数学模型的复杂性、手头的应用、可用资源以及许多其他标准。

我们在第2.2节开始本章，回顾线性的连续时间和离散时间系统的状态空间描述。在第2.3节中介绍了非线性系统的类似回顾。非线性系统的平衡在第2.4节中回顾，而向量空间的概念在第2.5节中介绍。稳定性理论在第2.6节中介绍，这是本章的主要内容。在第2.7节中介绍了李雅普诺夫稳定性结果，而输入-输出稳定性概念在第2.8节中介绍。在第2.9节中编制了高级稳定性概念，以使后续发展更加简洁。在第2.10节中，我们回顾了一些有用的定理和引理。在第2.11节中，我们从状态空间的角度回顾了基本的线性控制器设计，本章在第2.12节中结束。

\section{线性状态变量系统}

本书中考虑的机器人等许多物理系统由微分或差分方程描述。这些描述方程通常从基本物理定律中获得，为系统的分析和控制提供了起点。当然，有些系统非常复杂，没有可用的描述微分（或差分）方程。我们在本书中不考虑这些系统。

在本节中，我们研究物理系统的状态空间模型，这些系统是线性的。我们将自己限制为由常微分方程描述的系统中，这将导致有限维状态空间。需要偏微分方程（导致无限维系统）来研究柔性机器人机械臂，但这些在本教材中不予考虑。我们强调，本章的材料旨在作为这些主题的快速介绍，并不全面。读者可参考[Kailath 1980]、[Antsaklis and Michel 1997]以了解线性控制系统的更严格介绍。

\subsection{连续时间系统}

在深入研究连续时间系统的数学描述之前，我们先明确几个重要的基本概念：

{\small\itshape
\hspace*{2em}\textbf{[校对者注：基本概念通俗解释]}
\begin{itemize}
    \item \textbf{连续时间 vs 离散时间}：连续时间系统的状态随时间\textit{连续变化}，在每一时刻都有定义（如真实的物理世界）；而离散时间系统只在特定的采样时刻有定义（如计算机程序）。机器人手臂的平滑运动是连续的，而数字控制器的输出是离散的。
    
    \item \textbf{动态 vs 非动态}：动态系统的当前输出\textit{取决于过去}的输入和状态（有"记忆"），如机器人关节的位置取决于之前的历史运动；非动态系统的输出只取决于当前输入（无记忆、即时响应），如纯比例控制器 $u(t) = K_p e(t)$。
    
    \item \textbf{时不变 vs 时变}：时不变系统的特性不随时间改变（"脾气"稳定），如理想的刚性机械臂；时变系统的参数会随时间变化（如齿轮磨损导致间隙增大、负载变化导致惯量改变）。
\end{itemize}
}

连续时间系统被称为线性的，如果它遵循叠加原理，即如果输出 $y_1(t)$ 由输入 $u_1(t)$ 产生，输出 $y_2(t)$ 由输入 $u_2(t)$ 产生，那么由 $u(t)=\alpha_1 u_1(t)+\alpha_2 u_2(t)$ 产生的输出由 $y(t)=\alpha_1 y_1(t)+\alpha_2 y_2(t)$ 给出，其中 $\alpha_1$ 和 $\alpha_2$ 是标量常数。线性、单输入/单输出（SISO）、连续时间、时不变系统由线性、标量、常系数常微分方程描述，例如：
\begin{equation}\label{eq:2.2.1}
\sum_{i=0}^{n} a_i \frac{d^i y(t)}{dt^i} = \sum_{i=0}^{n} b_i \frac{d^i u(t)}{dt^i}
\end{equation}
其中 $a_i, b_i, i=0,\ldots,n$ 是标量常数，$y(t)$ 是标量输出，$u(t)$ 是标量输入。此外，在时间 $t_0$ 给定初始条件：
\begin{equation}
y(t_0), \frac{dy(t_0)}{dt}, \ldots, \frac{d^{n-1}y(t_0)}{dt^{n-1}}
\end{equation}
注意，输入 $u(t)$ 的导数次数最多与输出 $y(t)$ 相同。否则，系统被称为非动态的。

\subsubsection{状态空间实现}

系统的状态被定义为一组充分的变量，当在时间 $t_0$ 指定这些变量以及输入 $u(t), t \geq t_0$ 时，足以完全确定系统在 $t \geq t_0$ 时的行为[Kailath 1980]。状态向量然后包含确定系统中任何信号的未来行为所需的所有必要变量。根据定义，这样的状态向量 $x(t)$ 不是唯一的，这一特性稍后将加以利用。事实上，如果 $x$ 是状态向量，那么任何 $\tilde{x}(t)=Tx(t)$ 也是状态向量，其中 $T$ 是任何 $n \times n$ 可逆矩阵。对于在\eqref{eq:2.2.1}中描述的连续时间系统，可以选择以下状态向量：
\begin{equation}\label{eq:2.2.2}
x_i(t) = \frac{d^{i-1}y(t)}{dt^{i-1}}, \quad i=1,2,\ldots,n
\end{equation}
其中 $x_1=y, x_2=\dot{y}, \ldots, x_n=y^{(n-1)}$。输入-输出方程然后简化为：
\begin{equation}\label{eq:2.2.3}
y(t) = b_0 x_1(t) + b_1 x_2(t) + \cdots + b_{n-1} x_n(t) + u(t)
\end{equation}

\eqref{eq:2.2.2}和\eqref{eq:2.2.3}的更紧凑表述由下式给出：
\begin{equation}\label{eq:2.2.4}
\begin{aligned}
\dot{x} &= Ax + Bu \\
y &= Cx + Du
\end{aligned}
\end{equation}
其中
\begin{equation}\label{eq:2.2.5}
A = \begin{bmatrix}
0 & 1 & 0 & \cdots & 0 \\
0 & 0 & 1 & \cdots & 0 \\
\vdots & \vdots & \vdots & \ddots & \vdots \\
0 & 0 & 0 & \cdots & 1 \\
-\frac{a_0}{a_n} & -\frac{a_1}{a_n} & -\frac{a_2}{a_n} & \cdots & -\frac{a_{n-1}}{a_n}
\end{bmatrix}, \quad
B = \begin{bmatrix} 0 \\ 0 \\ \vdots \\ 0 \\ 1 \end{bmatrix}
\end{equation}

这种特定的状态空间表示被称为可控规范型[Kailath 1980]、[Antsaklis and Michel 1997]。一般来说，线性、时不变、连续时间系统将具有多个输入和一个以上的输出。事实上，$u(t)$ 是 $m \times 1$ 向量，$y(t)$ 是 $p \times 1$ 向量。将 $u(t)$ 与 $y(t)$ 关联的微分方程此处不介绍，但多输入/多输出（MIMO）系统的状态空间表示变为：
\begin{equation}\label{eq:2.2.6}
\begin{aligned}
\dot{x} &= Ax + Bu \\
y &= Cx + Du
\end{aligned}
\end{equation}
其中 $A$ 是 $n \times n$，$B$ 是 $n \times m$，$C$ 是 $p \times n$，$D$ 是 $p \times m$。有关 $A, B, C$ 和 $D$ 的具体形式，读者再次参考[Kailath 1980]、[Antsaklis and Michel 1997]。\eqref{eq:2.2.6}的框图如图2.2.1a所示。注意，状态的最小数量等于找到微分方程组的唯一解所需的初始条件的数量。

\begin{example}[双重积分器]
考虑由以下方程描述的SISO系统：
\[
\ddot{y}(t) = u(t)
\]
这个系统被称为双重积分器，代表了由牛顿定律描述的各种各样的物理系统。为了获得状态空间描述，设：
\[
x_1 = y, \quad x_2 = \dot{y}
\]
所以
\[
\dot{x}_1 = x_2, \quad \dot{x}_2 = u, \quad y = x_1
\]
\end{example}

\begin{example}[双平台系统]
考虑图2.2.2所示的MIMO机械系统，它代表了一个用于将实验与外部扰动隔离的双平台系统。系统有两个输入：导致地面移动的 $u_2$ 和导致平台 $m_1$ 移动的 $u_1$。系统还有2个输出，即平台 $m_1$ 的运动 $y_1$ 和平台 $m_2$ 的运动 $y_2$。实验将在平台 $m_1$ 的顶部进行，因此，人们希望最小化 $y_1$ 的大小。使用牛顿第二定律获得描述该系统的微分方程：
\[
m_1 \ddot{y}_1 + k_1(y_1 - y_2) + u_1 = 0
\]
\[
m_2 \ddot{y}_2 + k_1(y_2 - y_1) + k_2(y_2 - u_2) = 0
\]
可以通过选择以下状态获得该系统的状态空间表述：
\[
x_1 = y_1, \quad x_2 = \dot{y}_1, \quad x_3 = y_2, \quad x_4 = \dot{y}_2
\]
\end{example}

\subsubsection{传递函数}

线性、时不变、连续时间系统的另一种等价表示由其传递函数给出，它将系统的输入 $u(t)$ 与其输出 $y(t)$ 在拉普拉斯变量 $s$ 中或在频域中关联起来。重要的是要注意，传递函数描述没有关于系统初始条件的信息，因此，除非所有初始条件都为零，否则它不会为特定输入提供唯一输出[Antsaklis and Michel 1997]、[Kailath 1980]。然而，传递函数形式在实践中很重要，因为许多工程师熟悉频域规范。此外，许多系统的识别可以在频域中有效地进行[Ljung 1999]。因此，人们必须能够在状态空间（或现代）描述和传递函数（或经典）描述之间进行转换。

让我们考虑由\eqref{eq:2.2.6}描述的系统并取其拉普拉斯变换：
\begin{equation}\label{eq:2.2.7}
\begin{aligned}
sX(s) &= AX(s) + BU(s) + x(0) \\
Y(s) &= CX(s) + DU(s)
\end{aligned}
\end{equation}
其中 $X(s), U(s)$ 和 $Y(s)$ 分别是 $x(t), u(t)$ 和 $y(t)$ 的拉普拉斯变换。初始状态向量为 $x(0)$。通过在\eqref{eq:2.2.7}中的两个方程之间消去 $X(s)$，我们找到以下关系：
\begin{equation}\label{eq:2.2.8}
Y(s) = [C(sI-A)^{-1}B + D]U(s) + C(sI-A)^{-1}x(0)
\end{equation}
如前所述，当 $x(0)=0$ 时，传递函数作为输入 $U(s)$ 和输出 $Y(s)$ 之间的关系获得，即：
\begin{equation}\label{eq:2.2.9}
Y(s) = [C(sI-A)^{-1}B + D]U(s)
\end{equation}
这个特定线性、时不变系统的传递函数由下式给出：
\begin{equation}\label{eq:2.2.10}
P(s) = C(sI-A)^{-1}B + D
\end{equation}
使得（见图2.2.1）：
\begin{equation}\label{eq:2.2.11}
Y(s) = P(s)U(s)
\end{equation}

\begin{example}[双重积分器的传递函数]
考虑示例2.2.1的系统。很容易看出传递函数是：
\[
P(s) = \frac{1}{s^2}
\]
\end{example}

\subsection{离散时间系统}

在离散时间情况下，使用差分方程描述系统如下：
\begin{equation}\label{eq:2.2.12}
\sum_{i=0}^{n} a_i y(k+i) = \sum_{i=0}^{n} b_i u(k+i)
\end{equation}
其中 $a_i, b_i, i=0,\ldots,n$ 是标量常数，$y(k)$ 是输出，$u(k)$ 是时间 $k$ 的输入。注意，时间 $k+n$ 的输出依赖于时间 $k+n$ 的输入，但不依赖于后面的输入；否则，系统将是非因果的。

\subsubsection{状态空间表示}

以与连续时间情况类似的方式，定义以下状态向量：
\begin{equation}\label{eq:2.2.13}
x_i(k) = y(k+i-1), \quad i=1,2,\ldots,n
\end{equation}
输入-输出方程然后简化为：
\begin{equation}\label{eq:2.2.14}
y(k) = b_0 x_1(k) + b_1 x_2(k) + \cdots + b_{n-1}x_n(k) + u(k)
\end{equation}
\eqref{eq:2.2.13}和\eqref{eq:2.2.14}的更紧凑表述由下式给出：
\begin{equation}\label{eq:2.2.15}
\begin{aligned}
x(k+1) &= Ax(k) + Bu(k) \\
y(k) &= Cx(k) + Du(k)
\end{aligned}
\end{equation}
其中
\begin{equation}\label{eq:2.2.16}
A = \begin{bmatrix}
0 & 1 & 0 & \cdots & 0 \\
0 & 0 & 1 & \cdots & 0 \\
\vdots & \vdots & \vdots & \ddots & \vdots \\
0 & 0 & 0 & \cdots & 1 \\
-\frac{a_0}{a_n} & -\frac{a_1}{a_n} & -\frac{a_2}{a_n} & \cdots & -\frac{a_{n-1}}{a_n}
\end{bmatrix}
\end{equation}
MIMO情况与连续时间情况类似，由下式给出：
\begin{equation}\label{eq:2.2.17}
\begin{aligned}
x(k+1) &= Ax(k) + Bu(k) \\
y(k) &= Cx(k) + Du(k)
\end{aligned}
\end{equation}
其中 $A$ 是 $n \times n$，$B$ 是 $n \times m$，$C$ 是 $p \times n$，$D$ 是 $p \times m$。

在许多实际情况下，例如在机器人的控制中，系统是连续时间系统，但控制器使用数字硬件实现。这将要求设计者在连续和离散时间系统之间进行转换。有许多不同的方法来"离散化"连续时间系统，其中一些在第3章中讨论。对这一控制问题的这一非常重要方面感兴趣的读者可参考[Åström and Wittenmark 1996]、[Franklin et al. 1997]。

\begin{example}[离散时间双重积分器]
回忆示例2.2.1中提出的双重积分器或牛顿系统模型。微分方程的一个离散时间版本由以下差分方程给出：
\[
y(k+2) - 2y(k+1) + y(k) = T^2 u(k)
\]
其中 $T$ 是以秒为单位的采样周期。如果我们选择 $x_1(k)=y(k)$ 和 $x_2(k)=x_1(k+1)$，我们得到状态空间描述：
\[
\begin{aligned}
x_1(k+1) &= x_2(k) \\
x_2(k+1) &= 2x_2(k) - x_1(k) + T^2 u(k) \\
y(k) &= x_1(k)
\end{aligned}
\]
\end{example}

\subsubsection{传递函数表示}

以与连续时间情况类似的方式，由\eqref{eq:2.2.17}给出的线性、时不变、离散时间系统可以在Z变换域中由其传递函数 $P(z)$ 描述，从输入 $U(z)$ 到输出 $Y(z)$，使得：
\[
Y(z) = P(z)U(z)
\]
其中
\[
P(z) = C(zI-A)^{-1}B + D
\]
注意，Z变换用于离散时间情况，而拉普拉斯变换用于连续时间情况。

\begin{example}[离散时间双重积分器的传递函数]
示例2.2.4的传递函数由下式给出：
\[
P(z) = \frac{T^2}{z^2 - 2z + 1}
\]
\end{example}

\section{非线性状态变量系统}

在许多情况下，潜在的物理行为可能无法使用线性状态变量方程来描述。这就是机器人机械臂的情况，其中不同连杆之间的相互作用由非线性微分方程描述，如第3章所示。状态变量表述仍然能够处理这些系统，而传递函数和频域方法则失效。在本节中，我们处理前一节的非线性变体，并强调[Khalil 2001]、[Vidyasagar 1992]和[Verhulst 1997]、[LaSale and Lefschetz 1961]、[Hahn 1967]中研究的非线性系统的经典方法。

\subsection{连续时间系统}

非线性、标量、连续时间、时不变系统由非线性、标量、常系数微分方程描述，例如：
\begin{equation}\label{eq:2.3.1}
f(y, \dot{y}, \ldots, y^{(n)}, u, \dot{u}, \ldots, u^{(m)}) = 0
\end{equation}
其中 $y(t)$ 是输出，$u(t)$ 是所考虑系统的输入。与线性情况一样，我们通过其分量定义状态向量 $x$ 如下：
\begin{equation}\label{eq:2.3.2}
x_1 = y, \quad x_2 = \dot{y}, \quad \ldots, \quad x_n = y^{(n-1)}
\end{equation}
输出方程然后简化为：
\begin{equation}\label{eq:2.3.3}
y(t) = x_1(t)
\end{equation}
\eqref{eq:2.3.2}和\eqref{eq:2.3.3}的更紧凑表述由下式给出：
\begin{equation}\label{eq:2.3.4}
\begin{aligned}
\dot{x} &= f(x, U) \\
y &= cx
\end{aligned}
\end{equation}
其中
\[
U(t) = [u(t) \quad \dot{u}(t) \quad \cdots \quad u^{(n-1)}(t)]^T
\]
和
\begin{equation}\label{eq:2.3.5}
c = [1 \quad 0 \quad 0 \quad \cdots \quad 0]
\end{equation}

\begin{example}[非线性系统]
我们提出2个示例来说明这些概念：
\begin{enumerate}
\item 考虑阻尼摆方程
\[
\ddot{y} + \dot{y} + \sin(y) = 0
\]
通过选择 $x_1=y, x_2=\dot{y}$ 获得状态空间描述，导致：
\[
\dot{x}_1 = x_2, \quad \dot{x}_2 = -x_2 - \sin(x_1)
\]
$y(t)$ 的时间历程如图2.3.1所示。

\item 一个经典的非线性系统是Van der Pol振荡器，其描述为：
\[
\ddot{y} + (y^2-1)\dot{y} + y = 0
\]
或在状态空间中，其中：
\[
\dot{x}_1 = x_2, \quad \dot{x}_2 = -(x_1^2-1)x_2 - x_1
\]
$y(t)$ 的时间历程和相平面图（即 $x_2$ 对 $x_1$）如图2.3.2所示。
\end{enumerate}
\end{example}

\begin{example}[刚性机器人动力学]
刚性机器人由以下方程描述：
\[
M(q)\ddot{q} + V_m(q, \dot{q})\dot{q} + G(q) = \tau
\]
其中 $M(q)$ 是 $n \times n$ 惯性矩阵，$q$ 及其导数是 $n \times 1$ 广义坐标向量，$V(q, \dot{q}), G(q)$ 和 $\tau$ 分别是包含速度相关力矩、重力力矩和输入力矩的 $n \times 1$ 向量。

在此示例中，我们将集中于将 $n$ 个耦合微分方程写入状态空间形式。事实上，设状态向量 $x$ 为：
\[
x = \begin{bmatrix} q \\ \dot{q} \end{bmatrix}
\]
输入向量为 $u=\tau$，假设输出向量为 $y=q$。由于刚性机器人的一些特殊性质（见第2章），矩阵 $M(q)$ 已知是可逆的，因此：
\[
\ddot{q} = M^{-1}(q)[\tau - V_m(q,\dot{q})\dot{q} - G(q)]
\]
或
\[
\dot{x} = \begin{bmatrix} \dot{q} \\ M^{-1}(q)[\tau - V_m(q,\dot{q})\dot{q} - G(q)] \end{bmatrix}
\]
其中
\[
f(x) = \begin{bmatrix} x_2 \\ M^{-1}(x_1)[-V_m(x_1,x_2)x_2 - G(x_1)] \end{bmatrix}, \quad g(x) = \begin{bmatrix} 0 \\ M^{-1}(x_1) \end{bmatrix}
\]
\end{example}

\subsection{离散时间系统}

非线性、标量、离散时间、时不变系统由非线性、标量、常系数差分方程描述，例如：
\begin{equation}\label{eq:2.3.6}
f(y(k), y(k+1), \ldots, y(k+n), u(k), u(k+1), \ldots, u(k+n)) = 0
\end{equation}
其中 $y(.)$ 和 $u$ 如前所定义。状态变量的简单选择将导致：
\begin{equation}\label{eq:2.3.7}
x_i(k) = y(k+i-1), \quad i=1,2,\ldots,n
\end{equation}
或更紧凑地：
\begin{equation}\label{eq:2.3.8}
\begin{aligned}
x(k+1) &= f(x(k), U(k)) \\
y(k) &= cx(k)
\end{aligned}
\end{equation}
其中 $U(k)$ 和 $c$ 的定义类似于\eqref{eq:2.3.4}中给出的那些。

\begin{example}[逻辑斯蒂方程]
考虑标量系统
\[
y(k+1) = y(k) + u(k)y(k)(1-y(k))
\]
这导致状态空间表示：
\[
x(k+1) = x(k) + u(k)x(k)(1-x(k)), \quad y(k) = x(k)
\]
\end{example}

我们不强调离散非线性系统的研究，因为机器人由微分方程描述。然而，如第4章所讨论的，机器人控制器通常使用数字控制器实现。因此，能够像[Åström and Wittenmark 1995]、[Franklin et al. 1997]中讨论的那样在连续和离散时间描述的非线性动态系统之间进行转换将是有利的。

\section{非线性系统与平衡点}

在本节中，我们集中于由\eqref{eq:2.2.15}描述的系统，附加要求是 $u(t)$ 被指定为状态 $x(t)$ 的函数，即：
\begin{equation}\label{eq:2.4.1}
\dot{x}(t) = f[t, x(t)], \quad x \in \mathbb{R}^n
\end{equation}
这将使我们能够集中于分析问题。我们需要几个定义，我们现在介绍它们。

\begin{definition}[自治系统]
如果 $f[t, x(t)]$ 不明确依赖于时间 $t$，即：
\[
\dot{x}(t) = f[x(t)]
\]
则系统\eqref{eq:2.4.1}被称为自治的。
\end{definition}

\begin{example}[非自治系统]
在示例2.3.1中介绍的2个系统都是自治的，而由以下方程描述的系统：
\[
\dot{x} = \frac{x}{t}
\]
不是自治的。
\end{example}

\begin{definition}[平衡点]
向量 $x_e \in \mathbb{R}^n$ 是\eqref{eq:2.4.1}在时刻 $t_0$ 的固定点或平衡点，如果：
\[
f[t_0, x_e] = 0
\]
\end{definition}

\begin{example}[自治系统的平衡点]
由以下方程描述的系统：
\[
\dot{x} = Ax
\]
是自治的，并且它在 $\mathbb{R}^n$ 的原点处有一个平衡点。
\end{example}

注意：
\begin{itemize}
\item 如果系统是自治的，那么时间 $t_0$ 处的平衡点也是所有其他时间的平衡点。
\item 如果 $x_e$ 是非自治系统\eqref{eq:2.4.1}在时间 $t_0$ 的平衡点，那么 $x_e$ 是\eqref{eq:2.4.1}对于所有 $t_1 \geq t_0$ 的平衡点。
\end{itemize}

\begin{example}[非自治系统的平衡点]
考虑系统
\[
\dot{x} = x - t
\]
这是非自治的。它没有平衡点。虽然看起来在时间 $t_0=1$ 它有2个平衡点 $x_{e1}=-1$ 和 $x_{e2}=1$。然而，这些不是时间 $t \geq 1$ 的平衡点，因为在时间 $t \geq 1$ 平衡的条件不成立。
\end{example}

一些书籍也使用术语"静止点"或"奇点"来表示平衡点。

\begin{example}[阻尼摆]
回忆示例2.3.1中的摆，让我们尝试找到它的平衡点。首先注意系统是自治的，因此我们不需要指定特定的时间 $t_0$，然后注意，当两者都满足时摆处于平衡：
\[
x_2 = 0 \quad \text{和} \quad \sin(x_1) = 0
\]
因此平衡点位于：
\[
x_e = \begin{bmatrix} n\pi \\ 0 \end{bmatrix}, \quad n = 0, \pm 1, \pm 2, \ldots
\]
显然，当摆以零速度垂直向上或垂直向下悬挂时，摆处于平衡。
\end{example}

\begin{definition}[孤立平衡点]
如果存在 $x_e$ 的邻域 $N$，除了 $x_e$ 本身外不包含其他平衡点，则\eqref{eq:2.4.1}在时间 $t_0$ 的平衡点 $x_e$ 被称为孤立的。
\end{definition}

\begin{example}[摆的平衡点]
摆的平衡点是孤立的。另一方面，由 $\dot{x}=0$ 描述的系统具有 $\mathbb{R}$ 中的任何点作为其平衡点，因此其平衡点都不是孤立的。
\end{example}

\section{向量空间、范数与内积}

在本节中，我们将讨论非线性微分方程及其解的一些性质。我们将需要许多概念，如向量空间和范数，我们将简要介绍它们。读者可参考[Boyd and Barratt]、[Desoer and Vidyasagar 1975]、[Khalil 2001]进行证明和详细说明。

\subsection{线性向量空间}

在大多数应用中，我们需要处理（线性）实向量和复向量空间，它们随后被定义。

\begin{definition}[线性向量空间]
一个实线性向量空间（相应地，复线性向量空间）是一个集合 $V$，配备2个二元运算：加法（+）和标量乘法（$\cdot$），使得：
\begin{enumerate}
\item $x+y \in V$，如果 $x \in V$ 和 $y \in V$
\item $x+y = y+x$
\item $(x+y)+z = x+(y+z)$
\item 存在零元素 $0 \in V$，使得 $x+0=x$，对所有 $x \in V$
\item 对每个 $x \in V$，存在唯一的 $y=-x$，使得 $x+y=0$
\item $\alpha \cdot x \in V$，如果 $x \in V$ 和 $\alpha$ 是实数（相应地，复数）
\item $\alpha \cdot (x+y) = \alpha \cdot x + \alpha \cdot y$
\item $(\alpha + \beta) \cdot x = \alpha \cdot x + \beta \cdot x$
\item $(\alpha \beta) \cdot x = \alpha \cdot (\beta \cdot x)$
\item 对每个 $x \in V$，我们有 $1 \cdot x = x$，其中 $1$ 是 $\mathbb{R}$（相应地，$\mathbb{C}$）中的单位元
\end{enumerate}
\end{definition}

\begin{example}[向量空间]
以下是具有相关标量域的线性向量空间：$\mathbb{R}^n$ 带有 $\mathbb{R}$，和 $\mathbb{C}^n$ 带有 $\mathbb{C}$。
\end{example}

\begin{definition}[子空间]
向量空间 $V$ 的子集 $M$ 如果它本身是线性向量空间，则它是一个子空间。$M$ 成为子空间的一个必要条件是它包含零向量。
\end{definition}

我们可以为向量空间配备许多函数。其中之一是内积，它将 $V$ 中的两个向量带到 $\mathbb{R}$ 或 $\mathbb{C}$ 中的标量，另一个是向量的范数，它将 $V$ 中的向量带到 $\mathbb{R}$ 中的正值。以下部分讨论向量的范数，然后是内积的部分。

\subsection{信号和系统的范数}

范数是距离和长度概念的推广。由于稳定性理论通常关注某些向量和矩阵的大小，我们在这里简要介绍本书中将使用的一些范数。我们将首先考虑定义在向量空间 $X$ 上的向量的范数，具有相关的实数标量域 $\mathbb{R}$，然后引入矩阵诱导范数、函数范数，最后是系统诱导范数或算子增益。

\subsubsection{向量范数}

我们通过回顾最熟悉的赋范空间开始讨论范数，即具有常数条目的向量空间。在下文中，$|a|$ 表示实数 $a$ 的绝对值，或者如果 $a$ 是复数则表示 $a$ 的幅值。

\begin{definition}[向量范数]
向量 $x$ 的范数 $||\cdot||$ 是定义在向量空间 $X$ 上的实值函数，使得：
\begin{enumerate}
\item $||x|| > 0$ 对所有 $x \in X$，当且仅当 $x=0$ 时 $||x||=0$
\item $||\alpha x|| = |\alpha| \cdot ||x||$ 对所有 $x \in X$ 和任何标量 $\alpha$
\item $||x+y|| \leq ||x|| + ||y||$ 对所有 $x,y \in X$
\end{enumerate}
\end{definition}

\begin{example}[向量范数(1)]
以下是 $X=\mathbb{R}^n$ 中的常见范数，其中 $\mathbb{R}^n$ 是具有实分量的 $n \times 1$ 向量的集合。
\begin{enumerate}
\item 1-范数：$||x||_1 = \sum_{i=1}^{n} |x_i|$
\item 2-范数：$||x||_2 = \sqrt{\sum_{i=1}^{n} x_i^2}$，也称为欧几里得范数
\item p-范数：$||x||_p = \left(\sum_{i=1}^{n} |x_i|^p\right)^{1/p}$
\item $\infty$-范数：$||x||_\infty = \max_i |x_i|$
\end{enumerate}
\end{example}

\begin{example}[向量范数(2)]
考虑向量
\[
x = \begin{bmatrix} -1 \\ 2 \end{bmatrix}
\]
那么，$||x||_1 = 3$，$||x||_2 = \sqrt{5}$，$||x||_\infty = 2$。
\end{example}

我们现在给出 $\mathbb{R}^n$ 中向量范数的一个重要性质，它将在后续中有用。

\begin{lemma}[范数等价性]
设 $||x||_a$ 和 $||x||_b$ 是向量 $x \in \mathbb{R}^n$ 的任意两个范数。那么存在有限的正常数 $k_1$ 和 $k_2$，使得：
\[
k_1 ||x||_a \leq ||x||_b \leq k_2 ||x||_a
\]
\end{lemma}
引理中的两个范数被称为等价的，这个特定性质对于 $\mathbb{R}^n$ 上的任何两个范数都成立。

\subsubsection{矩阵范数}

在系统应用中，特定向量 $x$ 可以由矩阵 $A$ 操作以获得另一个向量 $y=Ax$。为了关联 $x$ 和 $Ax$ 的大小，我们定义诱导矩阵范数如下。

\begin{definition}[诱导矩阵范数]
设 $||x||$ 是给定的 $x \in \mathbb{R}^n$ 的范数。那么每个 $m \times n$ 矩阵 $A$ 都有一个诱导范数，由下式定义：
\[
||A||_i = \sup_{x \neq 0} \frac{||Ax||}{||x||} = \sup_{||x||=1} ||Ax||
\]
其中 $\sup$ 表示上确界。
\end{definition}

检查所提出的范数是否满足定义2.5.3的条件总是必要的。新定义的矩阵范数也可以证明满足：
\[
||AB||_i \leq ||A||_i ||B||_i
\]
对所有 $n \times m$ 矩阵 $A$ 和所有 $m \times p$ 矩阵 $B$。

\subsubsection{函数范数}

接下来，我们回顾时间相关函数和函数向量的范数。这些构成一类重要的信号，将在控制机器人时遇到。

\begin{definition}[函数范数]
设 $f(.): [0, \infty) \rightarrow \mathbb{R}$ 是一致连续函数。函数 $f$ 如果对于任何 $\epsilon > 0$，存在 $\delta(\epsilon)$ 使得当 $|t_1-t_2|<\delta$ 时 $|f(t_1)-f(t_2)|<\epsilon$，则称 $f$ 是一致连续的。

那么，如果对于 $p \in [1, \infty)$：
\[
\int_0^{\infty} |f(t)|^p dt < \infty
\]
$f$ 被称为属于 $L_p$。

如果 $f$ 是有界的，则称 $f$ 属于 $L_\infty$，即如果：
\[
\sup_{t \geq 0} |f(t)| < \infty
\]
其中 $\sup f(t)$ 表示 $f(t)$ 的上确界，即大于或等于 $f(t)$ 最大值的最小数。$L_1$ 表示具有有限绝对面积的信号集合，而 $L_2$ 表示具有有限总能量的信号集合。
\end{definition}

\subsubsection{系统范数}

接下来我们想研究多输入多输出（MIMO）系统对多维信号的影响。换句话说，时变向量 $u(t)$ 通过MIMO系统 $H$ 时会发生什么？设 $H$ 是一个具有 $m$ 输入和 $l$ 输出的系统，因此其对输入 $u(t)$ 的输出由下式给出：
\[
y(t) = (Hu)(t)
\]
如果 $u$ 属于 $L_p$ 时 $Hu$ 属于 $L_p$，并且存在有限常数 $\gamma > 0$ 和 $b$ 使得：
\[
||y||_p \leq \gamma ||u||_p + b
\]
则我们说 $H$ 是 $L_p$ 稳定的。

如果 $p=\infty$，系统被称为有界输入有界输出（BIBO）稳定。

\begin{definition}[$L_p$增益]
系统 $H$ 的 $L_p$ 增益记为 $\gamma_p(H)$，是使得存在有限 $b$ 来验证该方程的最小 $\gamma$。
\end{definition}

因此，增益 $\gamma_p$ 表征了输入信号通过系统时的放大。以下引理表征了线性系统的增益，可在[Boyd and Barratt]中找到。

\begin{lemma}[线性系统的增益]
给定线性系统 $H$，使得输入 $u(t)$ 产生输出 $y(t)=(h*u)(t)=\int_0^t h(t-\tau)u(\tau)d\tau$，并假设 $H$ 是BIBO稳定的，那么：
\begin{enumerate}
\item $\gamma_\infty(H) = \int_0^\infty |h(t)| dt$
\item $\gamma_2(H) = \sup_{\omega} |H(j\omega)|$
\item $\gamma_1(H)$ 不提供简单的解析解，但可以通过线性规划有效地计算
\end{enumerate}
\end{lemma}

\subsection{内积}

内积是向量空间中两个向量之间的运算，它将使我们能够定义几何概念，如正交性和傅里叶级数等。以下定义了内积。

\begin{definition}[内积]
在向量空间 $V$ 上定义的内积是一个函数 $\langle \cdot, \cdot \rangle$，从 $V$ 到 $F$，其中 $F$ 是 $\mathbb{R}$ 或 $\mathbb{C}$，使得对所有 $x, y, z \in V$：
\begin{enumerate}
\item $\langle x, y \rangle = \langle y, x \rangle^*$，其中 $\langle \cdot, \cdot \rangle^*$ 表示复共轭
\item $\langle x, y+z \rangle = \langle x, y \rangle + \langle x, z \rangle$
\item $\langle \alpha x, y \rangle = \alpha \langle x, y \rangle$，对所有标量 $\alpha$
\item $\langle x, x \rangle \geq 0$，其中仅当 $x=0_V$ 时取等号
\end{enumerate}
\end{definition}

我们可以为任何内积定义一个范数：
\[
||x|| = \sqrt{\langle x, x \rangle}
\]
因此，范数是一个更一般的概念：一个向量空间可以有一个与之相关的范数，但没有内积。然而，反过来则不然。

\subsection{矩阵性质}

一些矩阵性质在研究动态系统的稳定性中起着重要作用。本书所需的性质收集在本节中。我们将假设读者熟悉基本的线性代数。

\begin{definition}[矩阵定性]
本定义中的所有矩阵都是实方阵。
\begin{itemize}
\item 正定：如果对所有 $x \in \mathbb{R}^n, x \neq 0$，有 $x^T A x > 0$，则实 $n \times n$ 矩阵 $A$ 是正定的。
\item 半正定：如果对所有 $x \in \mathbb{R}^n$，有 $x^T A x \geq 0$，则实 $n \times n$ 矩阵 $A$ 是半正定的。
\item 负定：如果对所有 $x \in \mathbb{R}^n, x \neq 0$，有 $x^T A x < 0$，则实 $n \times n$ 矩阵 $A$ 是负定的。
\item 半负定：如果对所有 $x \in \mathbb{R}^n$，有 $x^T A x \leq 0$，则实 $n \times n$ 矩阵 $A$ 是半负定的。
\item 不定：如果对于某些 $x \in \mathbb{R}^n$，有 $x^T A x > 0$，而对于其他 $x \in \mathbb{R}^n$，有 $x^T A x < 0$，则 $A$ 是不定的。
\end{itemize}
\end{definition}

\begin{theorem}[特征值测试]
设 $A=[a_{ij}]$ 是对称的 $n \times n$ 实矩阵。因此，$A$ 的所有特征值都是实数。那么我们有以下结论：
\begin{itemize}
\item 正定：如果所有特征值都是正的，则实 $n \times n$ 矩阵 $A$ 是正定的。
\item 半正定：如果所有特征值都是非负的，则实 $n \times n$ 矩阵 $A$ 是半正定的。
\item 负定：如果所有特征值都是负的，则实 $n \times n$ 矩阵 $A$ 是负定的。
\item 半负定：如果所有特征值都是非正的，则实 $n \times n$ 矩阵 $A$ 是半负定的。
\item 不定：如果某些特征值是正的而某些是负的，则实 $n \times n$ 矩阵 $A$ 是不定的。
\end{itemize}
\end{theorem}

\begin{theorem}[Rayleigh-Ritz]
设 $A$ 是实对称的 $n \times n$ 正定矩阵。设 $\lambda_{\min}$ 是最小特征值，$\lambda_{\max}$ 是 $A$ 的最大特征值。那么，对任何 $x \in \mathbb{R}^n$：
\[
\lambda_{\min} x^T x \leq x^T A x \leq \lambda_{\max} x^T x
\]
\end{theorem}

\begin{theorem}[Gerschgorin圆盘定理]
设 $A=[a_{ij}]$ 是对称的 $n \times n$ 实矩阵。假设：
\[
r_i = \sum_{j \neq i} |a_{ij}|
\]
如果所有对角元素都是正的，即 $a_{ii} > 0$，则矩阵 $A$ 是正定的。
\end{theorem}

\section{稳定性理论}

我们研究的第一个稳定性概念涉及自由系统的行为，或等效地，具有给定输入的强制系统的行为。换句话说，我们研究系统平衡点对系统初始条件变化的稳定性。然而，在这样做之前，我们回顾一些基本定义。这些定义将以连续非线性系统的形式陈述，理解离散非线性系统允许类似结果，而线性系统只是非线性系统的特例。

设 $x_e$ 是自由连续时间（可能是时变）非线性系统的平衡点（或固定）状态：
\begin{equation}\label{eq:2.6.1}
\dot{x} = f(x, t), \quad x \in \mathbb{R}^n
\end{equation}
即 $f(x_e, t) = 0$，其中 $x, f$ 是 $n \times 1$ 向量。

我们将首先回顾平衡点 $x_e$ 的稳定性，理解 $x(t)$ 状态的稳定性总是可以通过稍后讨论的变量平移来获得。我们使用的稳定性定义可在[Khalil 2001]、[Vidyasagar 1992]中找到。

\begin{definition}[稳定性定义]\label{def:stability}
在本定义的所有部分中，$x_e$ 是时间 $t_0$ 处的平衡点，$||\cdot||$ 表示先前定义的任何函数范数。
\begin{enumerate}
\item 稳定性：如果在时刻 $t_0$ 足够接近 $x_e$ 开始的状态将始终在以后的时间保持在 $x_e$ 附近，则 $x_e$ 在李雅普诺夫意义下是稳定的（SL）。更精确地说，如果对任何给定的 $\epsilon > 0$，存在正数 $\delta(\epsilon, t_0)$，使得如果：
\[
||x_0 - x_e|| < \delta
\]
则：
\[
||x(t) - x_e|| < \epsilon, \quad \forall t \geq t_0
\]
$x_e$ 在李雅普诺夫意义下是稳定的，如果它对于任何给定的 $t_0$ 都是稳定的。

\item 不稳定性：如果无论状态多接近 $x_e$ 开始，它都不会在以后的时间被限制在 $x_e$ 的附近，则 $x_e$ 在李雅普诺夫意义下是不稳定的（UL）。换句话说，如果它在 $t_0$ 不稳定，则 $x_e$ 是不稳定的。

\item 收敛性：如果从接近 $x_e$ 开始的状态最终收敛到 $x_e$，则 $x_e$ 在时间 $t_0$ 是收敛的（C）。换句话说，如果对任何正数 $\epsilon$，存在正数 $\delta_1(t_0)$ 和正数 $T(\epsilon_1, x_0, t_0)$，使得如果：
\[
||x_0 - x_e|| < \delta_1
\]
则：
\[
||x(t) - x_e|| < \epsilon, \quad \forall t \geq t_0 + T
\]
$x_e$ 是收敛的，如果它对于任何 $t_0$ 都是收敛的。

\item 渐近稳定性：如果从充分接近 $x_e$ 开始的状态将保持在附近并最终收敛到它，则 $x_e$ 在时间 $t_0$ 是渐近稳定的（AS）。更准确地说，如果 $x_e$ 在时间 $t_0$ 既是收敛的又是稳定的，则它在时间 $t_0$ 是渐近稳定的。如果它对任何 $t_0$ 都是渐近稳定的，则 $x_e$ 是渐近稳定的。

\item 全局渐近稳定性：如果任何初始状态将保持在 $x_e$ 附近并最终收敛到它，则 $x_e$ 是全局渐近稳定的（GAS）。换句话说，如果它在 $t_0$ 是稳定的，并且如果每个 $x(t)$ 随着时间趋于无穷大收敛到 $x_e$，则 $x_e$ 是全局渐近稳定的。如果它对任何 $t_0$ 都是全局渐近稳定的，则 $x_e$ 是全局渐近稳定的，并且系统在这种情况下被称为全局渐近稳定的，因为它只能有一个平衡点 $x_e$。
\end{enumerate}
\end{definition}

注意，稳定性和渐近稳定性是局部概念，因为如果初始扰动 $\delta$ 太大，随后的状态 $x(t)$ 可能会任意远离 $x_e$。因此，存在以 $x_e$ 为中心的区域：
\[
||x - x_e|| < r
\]
使得对从该区域开始的任何状态都会产生稳定性和渐近稳定性，但对从该区域外开始的状态则不会。该区域被称为 $x_e$ 的吸引域。如果吸引域是 $\mathbb{R}^n$，则平衡点 $x_e$ 是全局渐近稳定的。

还应注意，所有先前的稳定性定义都依赖于初始时间 $t_0$，因此吸引域可能随变化的初始时间而变化。如果系统\eqref{eq:2.6.1}独立于时间（或自治），则定义\ref{def:stability}中的稳定性概念实际上独立于 $t_0$，并且它们将等同于接下来定义的稳定性概念。另一方面，即使系统\eqref{eq:2.6.1}是时间相关的，我们也希望其稳定性性质不依赖于 $t_0$，因为这将为我们提供所需的鲁棒性程度。这使我们定义均匀稳定性概念[Khalil 2001]。

\begin{definition}[均匀稳定性]
在本定义的所有部分中，$x_e$ 是时间 $t_0$ 处的平衡点。
\begin{enumerate}
\item 均匀稳定性：如果定义\ref{def:stability}中的 $\delta(\epsilon, t_0)$ 独立于 $t_0$，则 $x_e$ 在 $[t_0, \infty)$ 上是均匀稳定的（US）。

\item 均匀收敛性：如果定义\ref{def:stability}中的 $\delta_1(t_0)$ 和 $T(\epsilon_1, x_0, t_0)$ 可以独立于 $t_0$ 选择，则 $x_e$ 在 $[t_0, \infty)$ 上是均匀收敛的（UC）。

\item 均匀渐近稳定性：如果 $x_e$ 既是均匀稳定的又是均匀收敛的，则它在 $[t_0, \infty)$ 上是均匀渐近稳定的（UAS）。

\item 全局均匀渐近稳定性：如果 $x_e$ 是均匀稳定的和均匀收敛的，则它是全局均匀渐近稳定的（GUAS）。

\item 全局指数稳定性：如果对所有 $x_0 \in \mathbb{R}^n$，存在 $\alpha > 0$ 和 $\beta \geq 0$ 使得：
\[
||x(t)|| \leq \beta ||x_0|| e^{-\alpha(t-t_0)}
\]
则 $x_e$ 是全局指数稳定的（GES）。
\end{enumerate}
\end{definition}

注意，GES意味着GUAS。

在许多情况下，只需要状态大小方面的界限。这是比李雅普诺夫稳定性更不严格的要求。研究下面有界性定义与定义\ref{def:stability}中李雅普诺夫稳定性定义之间的细微差别是有益的。

\begin{definition}[有界性]
\begin{enumerate}
\item 有界性：如果从接近 $x_e$ 开始的状态永远不会离得太远，则 $x_e$ 在时间 $t_0$ 是有界的（B）。换句话说，如果对于每个 $\delta > 0$，使得 $||x_0 - x_e|| < \delta$，存在正数 $\epsilon(r, t_0) < \infty$，使得对所有 $t \geq t_0$：
\[
||x(t) - x_e|| < \epsilon
\]
如果它对于任何 $t_0$ 都是有界的，则 $x_e$ 是有界的。

\item 均匀有界性：如果 $\epsilon(r, t_0)$ 可以独立于 $t_0$，则 $x_e$ 在 $[t_0, \infty)$ 上是均匀有界的（UB）。

\item 均匀最终有界性：如果从接近 $x_e$ 开始的状态最终会变得有界，则 $x_e$ 被称为均匀最终有界的（UUB）。更准确地说，如果对任何 $\delta > 0, \epsilon > 0$，存在有限时间 $T(\epsilon, \delta)$，使得每当 $||x_0 - x_e|| < \delta$ 时，以下满足：
\[
||x(t) - x_e|| < \epsilon, \quad \forall t \geq T(\epsilon, \delta)
\]

\item 全局均匀最终有界性：如果对 $\epsilon > 0$，存在有限时间 $T(\epsilon)$ 使得：
\[
||x(t) - x_e|| < \epsilon, \quad \forall t \geq T(\epsilon)
\]
则 $x_e$ 被称为全局均匀最终有界的（GUUB）。
\end{enumerate}
\end{definition}

\section{李雅普诺夫稳定性定理}

李雅普诺夫稳定性理论涉及由微分方程描述的非强制非线性系统的行为：
\begin{equation}\label{eq:2.7.1}
\dot{x} = f(t, x), \quad x \in \mathbb{R}^n
\end{equation}
其中不失一般性，原点是\eqref{eq:2.7.1}的平衡点。读者可能认为这样的理论是不必要的，因为我们在前一节的示例中所要做的就是求解微分方程，并研究状态向量范数的时间演化。至少有两个原因说明为什么需要李雅普诺夫理论。第一个是李雅普诺夫理论将允许我们确定特定平衡点的稳定性，而无需实际求解微分方程。这，如任何非线性微分方程的学生所熟知的那样，是一个巨大的节省。第二个相关的原因是，李雅普诺夫理论为我们提供了稳定性问题的定性结果，可用于设计非线性动力系统的稳定控制器。

我们将首先假设满足\eqref{eq:2.7.1}具有唯一解的任何必要条件[Khalil 2001]、[Vidyasagar 1992]。对应于 $x(t_0)=x_0$ 的唯一解是 $x(t, t_0, x_0)$，将简单地记为 $x(t)$。在我们实际引入李雅普诺夫定理之前，我们回顾某些函数类，这些函数将简化李雅普诺夫定理的陈述。

\subsection{K类函数}

考虑连续函数 $\alpha: \mathbb{R} \rightarrow \mathbb{R}$。

\begin{definition}[K类函数]
如果满足以下条件，我们说 $\alpha$ 属于K类：
\begin{enumerate}
\item $\alpha(0) = 0$
\item 对所有 $x > 0$，$\alpha(x) > 0$
\item $\alpha$ 是非递减的，即对所有 $x_1 > x_2$，有 $\alpha(x_1) \geq \alpha(x_2)$
\end{enumerate}
\end{definition}

\begin{example}[K类函数]
函数 $\alpha(x)=x^2$ 是一个K类函数。函数 $\alpha(x)=x^2+1$ 不是K类函数，因为条件(1)不成立。另一方面，$\alpha(x)=-x^2$ 不是K类函数，因为条件(2)和(3)不成立。
\end{example}

\begin{definition}[正定函数]
在以下内容中，$\mathbb{R}_+ = [0, \infty)$。
\begin{enumerate}
\item 局部正定：如果存在K类函数 $\alpha(.)$ 和 $\mathbb{R}^n$ 原点的邻域 $N$ 使得对所有 $t \geq 0$ 和所有 $x \in N$：
\[
V(t, x) \geq \alpha(||x||)
\]
则连续函数 $V: \mathbb{R}_+ \times \mathbb{R}^n \rightarrow \mathbb{R}$ 是局部正定的（l.p.d）。

\item 正定：如果 $N=\mathbb{R}^n$，则函数 $V$ 被称为正定的（p.d）。

\item 负定和局部负定：如果 $-V$ 是（局部）正定的，我们说 $V$ 是（局部）负定的（n.d）。
\end{enumerate}
\end{definition}

\begin{definition}[递减函数]
在以下内容中：
\begin{enumerate}
\item 局部递减：如果存在K类函数 $\beta(.)$ 和 $\mathbb{R}^n$ 原点的邻域 $N$ 使得对所有 $t \geq 0$ 和所有 $x \in N$：
\[
V(t, x) \leq \beta(||x||)
\]
则连续函数 $V: \mathbb{R}_+ \times \mathbb{R}^n \rightarrow \mathbb{R}$ 是局部递减的。

\item 递减：如果 $N=\mathbb{R}^n$，我们说 $V$ 是递减的。
\end{enumerate}
\end{definition}

\begin{definition}[沿轨迹的导数]
给定连续可微函数 $V: \mathbb{R}_+ \times \mathbb{R}^n \rightarrow \mathbb{R}$ 以及微分方程组\eqref{eq:2.7.1}，$V$ 沿\eqref{eq:2.7.1}的导数被定义为函数 $\dot{V}: \mathbb{R}_+ \times \mathbb{R}^n \rightarrow \mathbb{R}$，由下式给出：
\[
\dot{V}(t, x) = \frac{\partial V(t, x)}{\partial t} + \frac{\partial V(t, x)}{\partial x} f(t, x)
\]
\end{definition}

\subsection{李雅普诺夫定理}

我们现在准备陈述李雅普诺夫定理，我们将它们分组在定理\ref{thm:lyapunov}中。证明见[Khalil 2001]、[Vidyasagar 1992]。

\begin{theorem}[李雅普诺夫稳定性定理]\label{thm:lyapunov}
给定非线性系统：
\[
\dot{x} = f(t, x), \quad x \in \mathbb{R}^n
\]
在原点处有平衡点，即 $f(t, 0)=0$，并设 $N$ 是大小为 $\epsilon$ 的原点的邻域，即：
\[
N = \{x \in \mathbb{R}^n; ||x|| \leq \epsilon\}
\]
那么：
\begin{enumerate}
\item 稳定性：如果存在标量函数 $V(t, x)$ 具有连续偏导数，使得对 $x \in N$：
\begin{enumerate}
\item $V(t, x)$ 是正定的
\item $\dot{V}$ 是半负定的
\end{enumerate}
则原点在李雅普诺夫意义下是稳定的。

\item 均匀稳定性：如果除了(a)和(b)之外，$V(t, x)$ 对 $x \in N$ 是递减的，则原点是均匀稳定的。

\item 渐近稳定性：如果 $V(t, x)$ 满足(a)并且对 $x \in N$ 是负定的，则原点是渐近稳定的。

\item 全局渐近稳定性：如果 $V(t, x)$ 验证(a)，并且对所有 $x \in \mathbb{R}^n$（即 $N=\mathbb{R}^n$）$\dot{V}(t, x)$ 是负定的，则原点是全局渐近稳定的。

\item 均匀渐近稳定性：如果 $V(t, x)$ 满足(a)，$V(t, x)$ 是递减的，并且 $\dot{V}(t, x)$ 对 $x \in N$ 是负定的，则原点是UAS。

\item 全局均匀渐近稳定性：如果 $N=\mathbb{R}^n$，并且如果 $V(t, x)$ 满足(a)，$V(t, x)$ 是递减的，$\dot{V}(t, x)$ 是负定的，并且 $V(t, x)$ 是径向无界的（即如果时间趋于一致，它趋于无穷大当 $||x|| \rightarrow \infty$），则原点是GUAS。

\item 指数稳定性：如果存在正常数 $\alpha, \beta, \gamma$ 使得：
\[
\alpha ||x||^2 \leq V(t, x) \leq \beta ||x||^2
\]
\[
\dot{V}(t, x) \leq -\gamma ||x||^2
\]
则原点是指数稳定的。

\item 全局指数稳定性：如果它对所有 $x \in \mathbb{R}^n$ 是指数稳定的，则原点是全局指数稳定的。
\end{enumerate}
\end{theorem}

定理中的函数 $V(t, x)$ 被称为李雅普诺夫函数。注意，该定理提供了原点稳定性的充分条件，无法为特定系统提供李雅普诺夫函数候选并不表示原点的稳定性。

\subsection{自治情况}

假设开环系统不是自治的，即不明确依赖于 $t$，那么可以获得与时间无关的李雅普诺夫函数候选 $V(x)$，并且正定条件大大简化，如下所述。

\begin{lemma}
如果 $V(0)=0$ 并且对所有 $x \neq 0$，$V(x) > 0$，则时间不变的连续函数 $V(x)$ 是正定的。如果上述在邻域中成立，则它是局部正定的。
\end{lemma}

注意，条件 $V(0)=0$ 不是必需的，并且只要 $V(0)$ 有上界，李雅普诺夫结果就无需修改而成立。

利用上述简化，李雅普诺夫结果成立，只是均匀和常规稳定性结果之间没有区别。

\begin{theorem}[LaSalle定理]
给定自治非线性系统：
\[
\dot{x} = f(x)
\]
并设原点是平衡点。那么：
\begin{enumerate}
\item 渐近稳定性：假设已经找到了一个李雅普诺夫函数 $V(x)$，使得对 $x \neq 0$，$V(x) > 0$ 并且 $\dot{V} \leq 0$。那么，当且仅当 $\dot{V}=0$ 仅在 $x=0$ 处时，原点是渐近稳定的。
\item 全局渐近稳定性：如果上述 $\dot{V}=0$ 并且 $V(x)$ 是径向无界的，则原点是GAS。
\end{enumerate}
\end{theorem}

\subsection{线性时不变情况}

在所考虑的系统是线性且时不变的情况下，李雅普诺夫理论发展得很好，并且李雅普诺夫函数的选择很简单。事实上，在这种情况下，定义2.6.1中的各种稳定性概念是相同的。李雅普诺夫理论然后提供了稳定性的必要和充分条件，如本节所讨论的。证明参考[Khalil 2001]。

\begin{theorem}
给定线性时不变系统：
\[
\dot{x} = Ax
\]
当且仅当存在方程的正定解 $P$ 时，系统是稳定的：
\[
A^T P + PA = -Q
\]
其中 $Q$ 是任意正定矩阵。
\end{theorem}

注意，在最后这个定理中获得了整个系统的稳定性，因为在这种情况下，原点是唯一的平衡点，其稳定性等同于系统的稳定性。此外，没有提到暗示什么类型的稳定性，因为在线性、时不变系统的非常特殊情况下，所有稳定性概念都是等价的[Khalil 2001]。还要注意，这个结果等同于测试 $A$ 的所有特征值是否具有负实部[Kailath 1980]。

\begin{theorem}[Krasovskii定理]
考虑自治非线性系统 $\dot{x}=f(x)$，原点是平衡点。设 $A(x)=\partial f/\partial x$。那么，原点为AS的一个充分条件是存在2个对称正定矩阵 $P$ 和 $Q$，使得对所有 $x \neq 0$，矩阵：
\[
F(x) = A(x)^T P + PA(x) + Q
\]
在某个以原点为中心的球 $B$ 中是正定的。函数 $V(x)=f(x)^T P f(x)$ 然后是系统的李雅普诺夫函数。如果 $B=\mathbb{R}^n$ 并且如果 $V(x)$ 是径向无界的，则系统是GAS。
\end{theorem}

\section{输入/输出稳定性}

当处理非线性系统时，李雅普诺夫意义下的稳定性并不一定意味着有界输入将导致有界输出。这个事实在下一个示例中显示。

然而，我们需要讨论有界输入将导致有界输出的条件[Boyd and Barratt]、[Desoer and Vidyasagar 1975]。这实际上在讨论系统诱导范数时已经提出（见定义2.5.7），当前的讨论应作为这些概念与李雅普诺夫稳定性的对比。

\begin{definition}[BIBO稳定性]
如果对于任何 $||u(t)|| \leq M < \infty$，存在有限 $\gamma > 0$ 和 $b$ 使得：
\[
||y(t)|| \leq \gamma M + b
\]
则动态系统\eqref{eq:2.8.1}-\eqref{eq:2.8.2}是有界输入有界输出（BIBO）稳定的。
\end{definition}

注意，BIBO稳定性意味着所有平衡状态的均匀有界性。

\section{高级稳定性结果}

在本节中，我们回顾一些高级稳定性概念。这些结果将用于显示使用鲁棒或自适应控制器时系统的闭环稳定性。如果读者只对实现这些控制器感兴趣，可以跳过本节。另一方面，任何对设计新控制器感兴趣的人都应该了解这里介绍的结果。

\subsection{被动系统}

给定图2.9.1所示的非线性系统，我们感兴趣的是仅基于输入-输出测量来研究这种系统的稳定性。受网络理论中能量概念的启发，例如：
\[
\text{[储存的能量]} = \text{[外部功率输入]} + \text{[内部功率产生]}
\]
人们可以研究所有类型系统的内部稳定性。一般来说，外部功率输入是输入努力或流 $u$ 与输出流或努力 $y$ 的标量积 $y^T u$。最后一个方程然后采取以下形式：
\begin{equation}\label{eq:2.9.1}
\dot{V}(x) = y^T u - g(t)
\end{equation}
在许多情况下（例如孤立系统），$g(t)=0$，可以使用储存的能量：
\begin{equation}\label{eq:2.9.2}
V(x) = \int_0^t y^T(\tau)u(\tau) d\tau
\end{equation}
作为李雅普诺夫函数候选。

非线性系统的被动性定义如下[Narendra and Taylor 1973]、[Ortega et al. 1998]。

\begin{definition}[被动性]
考虑图2.9.1所示的系统，并假设它具有相同数量的输入和输出，即 $u(t)$ 和 $y(t)$ 具有相同的维数。
\begin{enumerate}
\item 被动性：如果对所有有限 $T > 0$ 和一些 $\gamma > -\infty$：
\[
\int_0^T y^T(t)u(t) dt \geq -\gamma
\]
则系统被称为被动的。

\item 严格被动性：如果存在 $\delta > 0$ 和 $\gamma > -\infty$ 使得对所有有限 $T > 0$：
\[
\int_0^T y^T(t)u(t) dt \geq \delta \int_0^T ||u(t)||^2 dt - \gamma
\]
则系统被称为严格被动的。
\end{enumerate}
\end{definition}

被动系统实际上是一个不产生能量的系统。

\subsection{正实系统}

如果所考虑的系统是线性且时不变的，则被动性等同于正实性，可以在频域中测试[Narendra and Taylor 1973]。事实上，让我们描述正实系统并讨论它们的一些性质。考虑多输入多输出线性时不变系统：
\[
\begin{aligned}
\dot{x} &= Ax + Bu \\
y &= Cx + Du
\end{aligned}
\]
其中 $x$ 是 $n$ 向量，$u$ 是 $m$ 向量，$y$ 是 $p$ 向量，$A, B, C$ 和 $D$ 具有适当的维数。相应的传递函数矩阵是：
\[
P(s) = C(sI-A)^{-1}B + D
\]

我们将假设系统具有相等数量的输入和输出，即 $p=m$。

\begin{definition}[正实]
如果满足以下条件，则 $m \times m$ 矩阵 $T(s)$ 的正则有理函数（不恒为零）是正实的或PR：
\begin{enumerate}
\item $T(s)$ 的所有元素在区域 $\text{Re}(s) > 0$ 中没有极点
\item $T(s)$ 在 $j\omega$ 轴上的任何极点都是简单的，具有正定留数
\item 矩阵 $\text{He}[T(s)]$ 对 $\text{Re}(s) > 0$ 是半正定的
\end{enumerate}
\end{definition}

\begin{definition}[严格正实]
如果满足以下条件，则 $m \times m$ 矩阵 $T(s)$ 的正则有理函数（不恒为零）是严格正实的或SPR：
\begin{enumerate}
\item $T(s)$ 的所有元素在区域 $\text{Re}(s) \geq 0$ 中没有极点
\item 矩阵 $\text{He}[T(s)]$ 对 $\text{Re}(s) > 0$ 是正定的
\end{enumerate}
\end{definition}

\subsection{Lure问题}

考虑以下反馈互连系统：
\begin{equation}\label{eq:2.9.3}
\begin{aligned}
\dot{x} &= Ax + b u \\
y &= c^T x + d u \\
u &= -\phi(t, y)
\end{aligned}
\end{equation}
其中 $\phi$ 在两个参数中都是连续的。Lure随后陈述了以下绝对稳定性问题，后来被称为Lure问题：假设上述方程描述的系统给定，其中：
\begin{enumerate}
\item $A$ 的所有特征值都具有负实部，或者 $A$ 在原点处有一个特征值，而其余的特征值在开左半平面（OLHP）中
\item $(A, b)$ 是可控的
\item $(c, A)$ 是可观测的
\item 非线性 $\phi(.,.)$ 满足：
\begin{enumerate}
\item $\phi(t, 0) = 0$；$t \geq 0$
\item $0 \leq y\phi(t, y) \leq k y^2$；$y \in \mathbb{R}$
\end{enumerate}
\end{enumerate}
那么，找到线性系统 $(A, B, C, D)$ 上的条件，使得 $x=0$ 是闭环系统的GAS平衡点。

\subsection{MKY引理}

以下引理是Meyer-Kalman-Yakubovich（MKY）引理的版本，出现在[Narendra and Taylor 1973]、[Khalil 2001]等地，在设计机器人的自适应控制器时将非常有用。

\begin{lemma}[Meyer-Kalman-Yakubovitch]
设系统\eqref{eq:2.9.3}且 $D=0$ 是可控的。那么传递函数 $c(sI-A)^{-1}b$ 是SPR当且仅当：
\begin{enumerate}
\item 对任何对称正定 $Q$，存在对称正定 $P$ 解李雅普诺夫方程：
\[
A^T P + PA = -Q
\]
\item 矩阵 $B$ 和 $C$ 满足：
\[
C = B^T P
\]
\end{enumerate}
\end{lemma}

MKY引理给出了传递矩阵具有一定鲁棒性的条件。注意，条件取决于输入和输出矩阵，因此特定系统可能对某些输入/输出对的选择是SPR，而对其他选择则不是SPR。

\section{有用的定理与引理}

考虑图2.9.2所示的框图。标记为 $H_1$ 和 $H_2$ 的块代表两个系统（线性或非线性），它们对输入 $e_1$ 和 $e_2$ 进行如下操作：
\[
y_1 = H_1 e_1 = H_1(u_1 - y_2)
\]
\[
y_2 = H_2 e_2 = H_2(u_2 + y_1)
\]

\subsection{小增益定理}

\begin{theorem}[小增益定理]
设 $H_1: L_{pe} \rightarrow L_{pe}$ 和 $H_2: L_{pe} \rightarrow L_{pe}$。因此 $H_1$ 和 $H_2$ 对所有 $T \in [0, \infty)$ 满足不等式：
\[
||y_1||_T \leq \gamma_1 ||e_1||_T + \beta_1
\]
\[
||y_2||_T \leq \gamma_2 ||e_2||_T + \beta_2
\]
并且假设 $\gamma_1 \gamma_2 < 1$。如果：
\[
u_1, \nu_2 \in L_p\]
则 $e_1, y_2 \in L_p$ 和 $y_1, e_2 \in L_p$。
\end{theorem}

基本上，小增益定理指出，如果环增益小于1，则两个系统的反馈互连是BIBO稳定的。换句话说，如果一个信号遍历反馈环并在幅度上减小，则闭环系统不会变得不稳定。

\subsection{总稳定性定理}

\begin{theorem}[总稳定性]
考虑由下式描述的状态空间系统：
\[
\dot{x} = Ax + f(t, x) + g(t, x)
\]
其中：
\begin{enumerate}
\item 系统 $\dot{x} = Ax$ 是指数稳定的，即存在某些 $a > 0$ 和 $K \geq 1$，使得 $||e^{At}|| \leq Ke^{-at}$
\item $f(t, 0) = 0$，即原点是 $f(t, x)$ 的平衡点
\item $||f(t, x_1) - f(t, x_2)|| \leq \beta_1 ||x_1 - x_2||$，对某些 $\beta_1 > 0$
\item $||g(t, x)|| \leq \beta_2 r$，对某些 $\beta_2 > 0$
\item $||g(t, x_1) - g(t, x_2)|| \leq \beta_2 ||x_1 - x_2||$
\item $||x_0|| \leq r/K$
\end{enumerate}
那么，存在\eqref{eq:2.10.1}的唯一解 $x(t)$ 并且：
\[
||x(t)|| \leq k_1 e^{-\alpha t} ||x(0)|| + k_2 r
\]
\end{theorem}

总稳定性定理将用于设计控制器，使系统的线性部分指数稳定。实际上，这个定理保证如果系统的线性部分"非常"稳定（指数稳定），则有界非线性的去稳定效应可能不足以使系统去稳定，状态将保持有界。

\subsection{Bellman-Gronwall引理}

\begin{lemma}[Bellman-Gronwall]
设 $x(.), a(.), b(.): [0, \infty) \rightarrow [0, \infty)$，且 $T \geq 0$。假设对所有 $t \in [0, T]$，以下不等式成立：
\[
x(t) \leq \int_0^t a(\tau)x(\tau)d\tau + b(t)
\]
那么对所有 $t \in [0, T]$：
\[
x(t) \leq \int_0^t a(\tau)b(\tau)\exp\left(\int_\tau^t a(\sigma)d\sigma\right)d\tau + b(t)
\]
如果 $b(t)$ 进一步是常数，则以下成立：
\[
x(t) \leq b \exp\left(\int_0^t a(\tau)d\tau\right)
\]
\end{lemma}

\subsection{Barbalat引理}

\begin{lemma}[Barbalat]
设 $f(t)$ 是 $t$ 的可微函数。
\begin{itemize}
\item 第一版本：如果 $\dot{f}(t)$ 是一致连续的并且 $\lim_{t \rightarrow \infty} f(t) = k < \infty$，则 $\lim_{t \rightarrow \infty} \dot{f}(t) = 0$。
\item 第二版本：如果 $f(t) \geq 0$，$\dot{f}(t) \leq 0$ 并且 $\ddot{f}$ 有界，则 $\lim_{t \rightarrow \infty} \dot{f}(t) = 0$。
\end{itemize}
\end{lemma}

\section{线性控制器设计}

控制设计的目的是使机器人以可预测和期望的方式响应一组输入信号。在本节中，我们回顾如何设计线性控制器，以使机器人-控制器组合的行为是可接受的。机器人控制器的第一个要求显然是其稳定性。因此，控制器的第一个功能是稳定机器人在空间中移动时的状态。事实上，我们现在可以提出我们希望解决的问题：

\textbf{线性控制设计问题}：考虑一个LTI系统。设计一个对输出 $y(t)$ 进行操作而不进行微分并生成输入 $u(t)$ 的反馈控制器，使得系统从任何初始状态 $x(0)$ 在有限时间内达到指定的期望最终状态 $x_d$。这个问题有2个部分：
\begin{enumerate}
\item 可控性：如果 $y=x$，问题可以解决吗？
\item 可观测性：如果是这样，我如何从 $y$ 得到 $x$？
\end{enumerate}

在下一节中，我们介绍这些概念并找到测试，使我们能够回答所提出的问题。

\begin{definition}[可控性]
如果存在输入 $u(t)$，对 $0 \leq t \leq t_1$，使得 $x(t_1)=0$，对某些有限 $t_1$，则状态 $x_0$ 是可控的。如果所有 $x_0 \in \mathbb{R}^n$ 都是可控的，则系统本身被称为可控的。
\end{definition}

\begin{theorem}[可控性测试]
LTI系统完全可控的必要和充分条件是 $n \times nm$ 矩阵：
\[
C = [B \quad AB \quad A^2B \quad \cdots \quad A^{n-1}B]
\]
是满秩的，即 $\text{rank} C = n$。
\end{theorem}

接下来，我们展示可控性足以通过放置 $A$ 矩阵的特征值来稳定系统\eqref{eq:2.11.1}。

\begin{theorem}
设 $(A, b)$ 是可控的。那么，存在常数增益矩阵 $K$，使得 $u=-Kx+v$ 将把 $A-bK$ 的特征值放置在 $s$ 平面中的任何位置。
\end{theorem}

在单输入情况下，可以使用Ackermann公式找到放置特征值所需的状态反馈增益（参见例如[Antsaklis and Michel 1997]）。假设期望的闭环特征值被指定为方程的根：
\begin{equation}\label{eq:2.11.3}
\phi(s) = s^n + d_{n-1}s^{n-1} + \cdots + d_1 s + d_0 = 0
\end{equation}
那么，状态反馈控制器由下式给出：
\begin{equation}\label{eq:2.11.4}
K = [0 \quad 0 \quad \cdots \quad 0 \quad 1]C^{-1}\phi(A)
\end{equation}
其中 $C$ 是\eqref{eq:2.11.1}的可控性矩阵，$\phi(A)$ 是通过在 $A$ 处求值 $\phi(s)$ 获得的矩阵，即：
\begin{equation}\label{eq:2.11.5}
\phi(A) = A^n + d_{n-1}A^{n-1} + \cdots + d_1 A + d_0 I
\end{equation}

\begin{definition}[可观测性]
如果知道 $u(t), y(t)$；$0 \leq t \leq t_1$ 足以唯一确定 $x_0$，则状态 $x(0)=x_0$ 被称为可观测的。注意一旦 $x(0)$ 已知，$t \geq 0$ 的 $x(t)$ 也就知道了。如果每个初始状态都是可观测的，则系统被称为完全可观测的。
\end{definition}

\begin{theorem}[可观测性测试]
LTI系统完全可观测的必要和充分条件是 $np \times n$ 矩阵：
\[
O = \begin{bmatrix} C \\ CA \\ CA^2 \\ \vdots \\ CA^{n-1} \end{bmatrix}
\]
是满秩的。
\end{theorem}

通过结合可观测性和可控性的概念，我们可以设计解决线性控制设计问题的补偿器。事实上，以下定理总结了线性控制设计。

\begin{theorem}
当且仅当状态空间实现既是可观测的又是可控的，线性控制设计问题对系统\eqref{eq:2.11.1}是可解的。
\end{theorem}

\begin{quote}
\textit{注：读者可以通过}\href{https://frankjie09.github.io/Robot-Manipulator-Control-CN/figures/PID_Tuning_Demo.html}{\textbf{交互式可视化工具}}\textit{实时调节PID参数，观察Kp、Ki、Kd对机器人关节跟踪性能的影响。}
\end{quote}

\section{总结与注记}

本章旨在回顾将在本书其余部分中使用的各种控制理论概念。重点是包含足够的材料，以便控制理论背景很少或没有的读者能够跟随机器人控制器的设计。经常进行简化以集中于本书中研究的系统，即机械系统和由拉格朗日-欧拉方程描述的系统。特别是，双重积分器系统被详细说明，因为它将来自在机械机械臂上应用初步非线性反馈。在第3章中，我们介绍了刚性机器人机械臂的动力学描述，在随后的章节中，本章回顾的控制概念将在这些机械臂上实现。

这里回顾的大多数结果的更多细节可在[Kailath 1980]、[Antsaklis and Michel 1997]关于线性系统情况的书籍中找到，以及在[Khalil 2001]、[Vidyasagar 1992]关于非线性系统情况的书籍中找到。

\section*{参考文献}
\addcontentsline{toc}{section}{参考文献}

\begin{enumerate}
\item [Anderson et al. 1986] B.D.O.Anderson and R.Bitmead and C.R.Johnson, Jr. and P.V.Kokotovic and R.L.Kosut and I.M.Y.Mareels and L.Praly and B.D.Riedle. ``Stability of Adaptive Systems: Passivity and Averaging Analysis''. MIT Press, Cambridge, MA, 1986.

\item [Antsaklis and Michel 1997] P.Antsaklis and A.N.Michel. ``Linear Systems''. McGraw Hill, New York, 1997.

\item [Åström and Wittenmark 1995] K.Åström and B.Wittenmark. ``Adaptive Control''. Addison-Wesley, Reading, MA, 2nd edition, 1995.

\item [Åström and Wittenmark 1996] K.Åström and B.Wittenmark. ``Computer-Controlled Systems: Theory and Design'', Prentice-Hall, Englewood Cliffs, NJ, 1996.

\item [Boyd and Barratt] S.Boyd and G.Barratt. ``Linear Controller Design: Limits of Performance'', Prentice-Hall, Englewood Cliffs, NJ, 1991.

\item [Desoer and Vidyasagar 1975] C.Desoer and M.Vidyasagar. ``Feedback Systems: Input-Output Properties'', Academic Press, New York, 1975.

\item [Franklin et al. 1997] G.Franklin and D.Powell and M.Workman. ``Digital Control of Dynamic Systems'', Addison-Wesley, Reading, MA, 1997.

\item [Gu 1988] G.Gu. ``Stabilizability Conditions for Multivariable Uncertain systems via Output Feedback'', IEEE Trans. Autom. Control, vol. AC-35, pp. 988--992, 1988.

\item [Hahn 1967] W.Hahn. ``Theory and Applications of Lyapunov's Direct Method''. Prentice-Hall, Englewood Cliffs, NJ, 1967.

\item [Horn and Johnson 1991] R.Horn and C.Johnson. ``Topics in Matrix Analysis''. Cambridge University Press, New York, 1991.

\item [Kailath 1980] T.Kailath. ``Linear Systems''. Prentice-Hall, Englewood Cliffs, NJ, 1980.

\item [Khalil 2001] H.Khalil. ``Nonlinear Systems''. Prentice-Hall, Englewood Cliffs, NJ, 2001.

\item [LaSale and Lefschetz 1961] P.LaSale and S.Lefschetz. ``Stability by Lyapunov's Direct Method''. Academic Press, New York, 1961.

\item [Ljung 1999] L.Ljung. ``System Identification-Theory for the User''. Prentice-Hall, Englewood Cliffs, NJ, 1999.

\item [Narendra and Taylor 1973] K.S.Narendra and J.H.Taylor. ``Frequency Domain Criteria for Absolute Stability''. Academic Press, New York, 1973.

\item [Ortega and Spong 1988] R.Ortega and M.Spong. ``Adaptive Motion Control of Rigid Robots: A Tutorial''. Proc. IEEE Conf. Dec. and Cont., pp. 1575--1584, 1988.

\item [Ortega et al. 1998] R.Ortega and A.Loria and P.J.Nicklasson and H.Sira-Ramirez. ``Passivity-based Control of Euler-Lagrange Systems''. Springer-Verlag, Berlin, 1998.

\item [Sastry and Bodson 1989] S.Sastry and M.Bodson. ``Adaptive Control: Stability, Convergence, and Robustness''. Prentice-Hall, Englewood Cliffs, NJ, 1989.

\item [Vidyasagar 1992] M.Vidyasagar. ``Nonlinear Systems Analysis''. Prentice-Hall, Englewood Cliffs, NJ, 1992.
\end{enumerate}

\ifstandalone
  \end{document}
\fi
